\chapter{Introduction}
\label{ch:intro}

\section{Radio frequency interference}
% - Radio telescopes record extremely weak signals from space. Human-made
%   transmitters (mobile networks, satellites, aircraft, broadcast, electrical
%   equipment) are millions of times stronger. This is RFI.
% - RFI corrupts the data. It has to be found and removed ("flagged") before
%   any science is done. Your Week-2 slides cover the RFI types: constant
%   narrowband lines, constant broadband bands, instantaneous wideband bursts,
%   patchy broadband blobs -- reuse that, with the figure.
% - Telescopes produce far too much data to flag by hand, so flagging must be
%   automatic.
% - Original motivation from your Week-2 deck: GMRT pulsar timing for InPTA,
%   where uncleaned RFI damages sub-microsecond timing precision.

\section{RFI flagging as image segmentation}
% - A spectrogram is a 2D image: one axis time, one axis frequency, each
%   pixel a measured power.
% - Flagging = deciding, for every pixel, "RFI or not". That is exactly the
%   task computer-vision people call semantic segmentation.
% - U-Net (Ronneberger et al. 2015) is the standard segmentation network.
%   Akeret et al. 2017 were the first to apply it to RFI: tf_unet.
% - Classical alternatives: thresholding, AOFlagger (Offringa et al. 2012),
%   the 2-sigma CRF method (Chen et al. 2023). Details in Chapter 2.

\section{Aims of this project}
% Say what you set out to do, in the order you actually did it:
%  1. Reproduce the authors' tf_unet on our own data.
%  2. Understand why it performed badly (F1 0.39) and fix it.
%  3. Build an improved model.
%  4. Test both on REAL telescope data with human ground truth (LOFAR),
%     not only on synthetic data.
%  5. Find out what actually makes the improved model better.

\section{Contributions}
% TODO -- DECIDE THIS AFTER THE NOVELTY DISCUSSION WITH YOUR SUPERVISOR.
% This section sets the argument of the whole report. Candidates, all backed
% by measurements:
%  - Diagnosis of tf_unet's failure: class weighting + the ReLU output layer
%    drive it into a dead state; per-image normalisation does the rest.
%    0.3879 -> 0.9317 without changing the model. [PART 1]
%  - A 593,842-parameter model scoring 0.6603 +/- 0.0040 max F1 on the 109
%    expert-labelled LOFAR baselines, above the best published result
%    (RFI-Net, 0.5979). [PART 13.11, PART 14]
%  - Evidence that the gain does NOT come from the added architectural
%    components (strip convolutions, ECA, residual blocks). [PART 18, 20]
%  - Dataset audits: train/test leakage in the published HERA split (31 of
%    140 test images) and in the LOFAR split (all 109 test images). [PART 2, 11.5]
%  - The synthetic-to-real gap: identical code drops 0.9317 -> 0.5482. [PART 12]
% Do NOT claim the strip-convolution idea as novel: prior art MARS,
% arXiv:2608.05546 [PART 4].

\section{Structure of this report}
% One line per chapter. Write last.
